\documentclass[a4paper,10pt]{report}\input{../0.Préambule/Préambule_cours_prof}\begin{document}

\newpage
\section*{Cylindre et volumes}
\addcontentsline{toc}{section}{Cylindre et volumes}

\begin{exer1}
Parmi les vues ci-dessous entoure celles qui peuvent être des vues de cylindres.

\begin{center}

\includegraphics[scale=0.5]{../40.Image/17. Exo-Cylindre et volume 1.png}
\end{center}
\end{exer1}

\begin{exer1}
Dessine à main levé un cylindre en perspective cavalière.
\end{exer1}

\begin{exer2}
Les deux figures ne sont pas en vraie grandeur. \\Quel chemin rouge est le plus long ?

\vspace{-2cm}
\begin{flushright}
\includegraphics[scale=0.3]{../40.Image/17. Exo-Cylindre et volume 2.png}
\end{flushright}
\end{exer2}

\vspace{-7mm}
\subsection*{Patron de cylindre}
\addcontentsline{toc}{subsection}{Patron de cylindre}

\vspace{-3mm}
\noindent \begin{minipage}{8cm}
\begin{exer1}
Indique sur le schéma à main levée les mesures correspondant à la vue en perspective de ce cylindre. (On fera les calculs sur le cahier si besoin)

\begin{center}
\includegraphics[scale=0.33]{../40.Image/17. Exo-Cylindre et volume 3.png}
\end{center}
\end{exer1}
\end{minipage}
\hspace{5mm}
\begin{minipage}{8cm}
\begin{exer1}
On considère le patron d'un cylindre de révolution. Complète le tableau. Si besoin, donne des valeurs arrondis au dixième.

\begin{center}
\renewcommand{\arraystretch}{1.5}
\begin{tabularx}{8cm}{|*{3}{>{\centering \arraybackslash}X|}}
\hline
\cellcolor{exer1!20}{Rayon du cercle de base} & \cellcolor{exer1!20}{Diamètre du cercle de base} & \cellcolor{exer1!20}{Longueur du rectangle} \\\hline
$\np{4}$cm & & \\\hline
 & $\np{6,2}$cm & \\\hline
 & & $\np{1,58}$cm \\\hline
\end{tabularx}
\end{center}
\end{exer1}
\end{minipage}

\begin{exer2}
Construis  un patron d'un cylindre de révolution de hauteur $3$cm ayant pour base un disque de rayon $1$cm.
\end{exer2}

\subsection*{Volume de cylindre}
\addcontentsline{toc}{subsection}{Volume de cylindre}

\begin{exer1}
Calcule le volume, en cm$^3$, d'un cylindre de $5$cm de rayon et de $10$cm de hauteur.
\end{exer1}

\begin{exer2}
Calcule le volume, en L, d'un cylindre de $2,5$dm de rayon et de $400$mm de hauteur.
\end{exer2}

\begin{exer3}
Pour fabriquer un puits dans son jardin,\\ Mme Poncet a besoin d'acheter du béton.

\medskip
\noindent A l'aide des caractéristiques du cylindre, détermine le \\ volume de  béton nécessaire à la construction de ce \\ puits arrondi au centimètre-cube près.

\vspace{-2.9cm}
\begin{flushright}
\includegraphics[scale=0.7]{../40.Image/17. Exo-Cylindre et volume 4.png}
\end{flushright}
\end{exer3}

\begin{exer3}
Ci-contre est représenté un cylindre dont la hauteur mesure $5$cm \\et le rayon du disque de base mesure $2$cm.

\noindent Une partie du cylindre est représentée grisée est formé par l'intersection de deux \\ demi plans passant par l'axe de révolution du cylindre et formant un angle de $100$\degre

\begin{enumerate}[font=\color{exer3}]
\item Déterminer le volume du cylindre
\item Par proportionnalité, déterminer le volume de la partie grisée.
\end{enumerate}

\vspace{-4cm}
\begin{flushright}
\includegraphics[scale=0.7]{../40.Image/17. Exo-Cylindre et volume 5.png}
\end{flushright}

\end{exer3}
\end{document}
